\chapter{Degradation-State Mappings}

Slow degradation is implemented through bounded damage states rather than by
directly forcing the J--V curve to decay. These states map to model parameters
such as trap density, contact resistance, shunt resistance, surface
recombination, and generation loss, allowing each mechanism to be tested before
the mechanisms are combined.

\section{Purpose of Damage-State Variables}

The COMSOL model uses damage-state variables to connect slow physical changes
to parameters that the Semiconductor Module and Electrical Circuit interface can
use. Each state is dimensionless and bounded. A value of 0 represents the
baseline condition for that mechanism, and a value of 1 represents the selected
maximum severity for the simulated study. The bounded form prevents a numerical
damage variable from growing without physical interpretation.

The model maps latent degradation states onto COMSOL-accessible parameters such
as trap density, ionic charge, series resistance, shunt resistance, surface
recombination, and photogeneration. These mappings are not a substitute for
calibration. They define testable hypotheses about which physical mechanism
produces which J--V signature.

\section{Trap Degradation State}

Trap growth is motivated by defect formation, interfacial disorder, ion-driven
chemical change, and absorber damage. The model can represent trap degradation
either as a spatial PDE $D_{\mathrm{trap}}(x,t)$ in the perovskite domain or as
a lumped global ODE when spatial resolution is not identifiable. The spatial
form is
\begin{equation}
\frac{\partial D_{\mathrm{trap}}}{\partial t}
= k_{\mathrm{trap}} f_{\mathrm{stress}} f_{\mathrm{ion}}
\left(1-D_{\mathrm{trap}}\right)
- k_{\mathrm{heal}} f_{\mathrm{rest}}D_{\mathrm{trap}}.
\label{eq:dtrap}
\end{equation}
The dependent variable is dimensionless, so the source terms must have units of
1/s. The rate constant $k_{\mathrm{trap}}$ has units of 1/s when
$f_{\mathrm{stress}}$ and $f_{\mathrm{ion}}$ are dimensionless multipliers.
The healing term is optional and is used only when recovery during rest is
being tested.

The trap state maps to an effective trap density through either a bounded
linear interpolation,
\begin{equation}
N_{t,\mathrm{eff}} =
N_{t0} + (N_{t,\max}-N_{t0})D_{\mathrm{trap}},
\label{eq:nt_linear}
\end{equation}
or a multiplicative form,
\begin{equation}
N_{t,\mathrm{eff}} =
N_{t0}\left(1+a_{Nt}D_{\mathrm{trap}}\right).
\label{eq:nt_mult}
\end{equation}
Both forms preserve units of m$^{-3}$ or cm$^{-3}$ depending on the COMSOL
parameter convention. The parameter $N_{t,\max}$, or equivalently $a_{Nt}$,
sets the physical severity. A saturated damage state at
$D_{\mathrm{trap}}=1$ therefore means that the selected maximum trap density
has been reached, not that the material contains every possible defect.

In the J--V response, increased $N_t$ shortens SRH lifetimes, increases
nonradiative recombination, lowers $V_{oc}$, and can reduce FF or $J_{sc}$ if
carrier collection becomes recombination-limited. The spatial form can localize
trap growth near interfaces where $u_{\mathrm{ion}}$ or
$\nabla u_{\mathrm{ion}}$ is large. The limitation is that the ion--trap
coupling function $f_{\mathrm{ion}}$ is not unique from J--V data alone and
requires calibration with ion-sensitive or defect-sensitive measurements.

\section{SRH Recombination Coupling}

The SRH recombination expression connects the trap mapping to electronic
transport:
\begin{equation}
R_{\mathrm{SRH}} =
\frac{np-n_i^2}
{\tau_p(n+n_1)+\tau_n(p+p_1)}.
\end{equation}
The lifetimes are
\begin{align}
\tau_n &= \frac{1}{\sigma_n v_{\mathrm{th}} N_{t,\mathrm{eff}}},\\
\tau_p &= \frac{1}{\sigma_p v_{\mathrm{th}} N_{t,\mathrm{eff}}}.
\end{align}
The units are consistent because $\sigma$ has units of m$^2$,
$v_{\mathrm{th}}$ has units of m/s, and $N_t$ has units of m$^{-3}$, so
$\sigma v_{\mathrm{th}}N_t$ has units of 1/s. As
$N_{t,\mathrm{eff}}$ increases, both lifetimes decrease and
$R_{\mathrm{SRH}}$ increases for a given carrier population. This is the
physical path from trap degradation to lower voltage and stronger nonradiative
loss.

\section{Contact Resistance}

Contact resistance represents voltage loss at the metal or transport-layer
contact. In this model it is associated with the rear Ag contact boundary or an
equivalent circuit element. The area-normalized form is
\begin{equation}
V_{\mathrm{terminal}} =
V_{\mathrm{device}} + J R_{s,A},
\label{eq:rs_area}
\end{equation}
where $J$ is current density and $R_{s,A}$ has units of $\Omega\,\mathrm{m^2}$.
If a total current $I$ and device area $A$ are used instead, the equivalent
relation is
\begin{equation}
V_{\mathrm{terminal}} =
V_{\mathrm{device}} + J A R_s.
\end{equation}

The degradation mapping is
\begin{equation}
R_{s,\mathrm{eff}} =
R_{s0}\left(1+a_{Rs}D_{Rs}\right).
\label{eq:rs_deg}
\end{equation}
Increasing $R_s$ primarily rounds the high-forward-bias J--V knee and lowers
fill factor. It has little direct effect on $V_{oc}$ because the terminal
current is near zero at open circuit. The COMSOL interpretation is lumped
contact or interfacial voltage loss, not a spatially resolved chemical reaction
at the Ag boundary.

If light-soaking or contact activation is included, conditioning can reduce the
effective series resistance:
\begin{equation}
R_{s,\mathrm{eff}} =
\frac{R_{s0}\left(1+a_{Rs}D_{Rs}\right)}
{1+a_{\mathrm{cond}}D_{\mathrm{cond}}}.
\label{eq:rs_cond}
\end{equation}
The denominator represents beneficial extraction or contact activation. This
term is proposed for future calibration and is not treated as validated in the
present results.

\section{Shunt Resistance}

Shunting is represented as a leakage pathway in parallel with the photovoltaic
device. In area-normalized form,
\begin{equation}
J_{\mathrm{sh}} = \frac{V}{R_{sh,A}},
\label{eq:jsh}
\end{equation}
where $R_{sh,A}$ has units of $\Omega\,\mathrm{m^2}$. Equivalently, the model
can use a shunt conductance
\begin{equation}
G_{sh,\mathrm{eff}}=
G_{sh0}+\Delta G_{sh}D_{Rsh}.
\end{equation}
The resistance mapping can be written as
\begin{equation}
R_{sh,\mathrm{eff}} =
\frac{R_{sh0}}{1+a_{Rsh}D_{Rsh}}.
\label{eq:rsh_deg}
\end{equation}

Decreasing $R_{sh}$ increases leakage current near short circuit and changes
the low-voltage slope of the J--V curve. The usual signatures are FF loss and,
when leakage is strong enough, reduced $V_{oc}$. Because this effect is
introduced through a circuit-equivalent pathway, it is interpreted as lumped
leakage rather than a spatially resolved pinhole or defect filament.

\section{Surface Recombination and Photogeneration Loss}

Surface recombination is represented by an interface recombination velocity
$S$, with units of m/s. Increasing $S$ raises interface recombination and can
lower $V_{oc}$ and FF when interface carrier concentrations are high. The
mapping has the same bounded structure:
\begin{equation}
S_{\mathrm{eff}}=S_0\left(1+a_S D_S\right).
\end{equation}
The limitation is observability: terminal J--V data may not separate surface
recombination from bulk trap recombination without additional observables.

Photogeneration loss can be represented by reducing the absorber generation
rate:
\begin{equation}
G_{\mathrm{eff}}=G_0\left(1-a_GD_G\right).
\end{equation}
This mapping tests absorber damage or optical loss. Its primary J--V signature
is reduced $J_{sc}$, with secondary effects on PCE and possibly FF. It should
not be used to represent recombination loss unless
independent optical or quantum-efficiency evidence supports that choice.

\section{Light Soaking and Self-Healing}

PSCs can show beneficial conditioning under illumination and partial recovery
after rest. Light soaking can improve extraction or contact behavior, often
appearing as higher FF and sometimes higher PCE. Self-healing can occur through
ion redistribution, trap relaxation, and contact or interfacial
re-equilibration after stress is removed.

A proposed conditioning state is
\begin{equation}
\frac{dD_{\mathrm{cond}}}{dt}
= k_{\mathrm{cond}}f_{\mathrm{light}}f_T
\left(1-D_{\mathrm{cond}}\right)
- k_{\mathrm{relax}}D_{\mathrm{cond}}.
\label{eq:dcond}
\end{equation}
The state is dimensionless, so $k_{\mathrm{cond}}$ and $k_{\mathrm{relax}}$
have units of 1/s when the multipliers are dimensionless. The proposed mapping
to series resistance is Eq.~\ref{eq:rs_cond}.

Optional trap healing can be represented by the negative term in
Eq.~\ref{eq:dtrap}:
\begin{equation}
\frac{dD_{\mathrm{trap}}}{dt}
= k_{\mathrm{trap}}f_{\mathrm{stress}}
\left(1-D_{\mathrm{trap}}\right)
- k_{\mathrm{heal}}f_{\mathrm{rest}}D_{\mathrm{trap}}.
\end{equation}
These terms are physically motivated, but they are not fully validated in the
present COMSOL model. They remain future implementation targets that require
controlled light-soaking and rest-recovery data.
